\documentclass[11pt]{article}
\usepackage[a4paper,margin=28mm]{geometry}
\usepackage[T1]{fontenc}
\usepackage[utf8]{inputenc}
\usepackage{lmodern}
\usepackage{microtype}
\usepackage{amsmath,amssymb,amsthm}
\usepackage{xcolor}
\usepackage{hyperref}
\usepackage{listings}
\usepackage{enumitem}

\hypersetup{colorlinks=true,linkcolor=blue!55!black,citecolor=blue!55!black,urlcolor=blue!55!black}
\newtheorem{theorem}{Theorem}[section]
\newtheorem{lemma}[theorem]{Lemma}
\newtheorem{proposition}[theorem]{Proposition}
\newtheorem{corollary}[theorem]{Corollary}
\theoremstyle{definition}
\newtheorem{definition}[theorem]{Definition}
\newtheorem{example}[theorem]{Example}
\theoremstyle{remark}
\newtheorem{remark}[theorem]{Remark}

\definecolor{leanblue}{RGB}{245,248,252}
\definecolor{leangreen}{RGB}{22,112,65}
\lstdefinelanguage{Lean}{
  morekeywords={abbrev,by,calc,cases,constructor,def,else,end,exact,example,
  have,if,import,induction,intro,let,namespace,obtain,rcases,refine,rfl,
  rw,simpa,simp,theorem,where,with},
  sensitive=true,
  morecomment=[l]{--},
  morecomment=[s]{/-}{-/},
  morestring=[b]"
}
\lstset{language=Lean,basicstyle=\ttfamily\small,backgroundcolor=\color{leanblue},
  commentstyle=\color{leangreen},keywordstyle=\color{blue!65!black}\bfseries,
  columns=fullflexible,keepspaces=true,breaklines=true,frame=single,
  framesep=5pt,xleftmargin=4pt,xrightmargin=4pt,showstringspaces=false,
  literate=
    {α}{{$\alpha$}}1
    {¬}{{$\neg$}}1
    {·}{{$\cdot$}}1
    {₁}{{$_1$}}1
    {₂}{{$_2$}}1
    {←}{{$\leftarrow$}}1
    {→}{{$\to$}}1
    {∀}{{$\forall$}}1
    {∃}{{$\exists$}}1
    {∣}{{$\mid$}}1
    {∧}{{$\land$}}1
    {∨}{{$\lor$}}1
    {≠}{{$\neq$}}1
    {≤}{{$\le$}}1
    {⊢}{{$\vdash$}}1
    {▸}{{$\triangleright$}}1
    {⟨}{{$\langle$}}1
    {⟩}{{$\rangle$}}1}
\newenvironment{leanbridge}{%
  \par\smallskip\noindent\textbf{Lean realization.}\par\smallskip%
}{%
  \par\smallskip%
}
\newcommand{\SeqA}{\mathcal{A}}
\newcommand{\revbin}{\operatorname{revbin}}

\title{A Certified Positivity Theorem for OEIS A348369\\
\large Classical proof and Lean formalization in parallel}
\author{Thomas Scheuerle}
\date{\today}

\begin{document}
\maketitle

\begin{abstract}
Let $\SeqA$ be the increasing sequence OEIS A328596, whose elements are the
positive integers whose reversed binary expansions are Lyndon words.  OEIS
A348369 counts, for the $n$th element of $\SeqA$, its unordered representations
as a sum of two earlier elements of $\SeqA$.  We prove that every element of
$\SeqA$ other than $1$ has such a representation.  Consequently,
$a(n)\geq 1$ for every $n>1$ in A348369.  The proof removes the
least-significant set bit.  Each mathematical definition and argument below
is immediately followed by the corresponding Lean 4 fragment, so the
ordinary proof and its machine-checked realization can be read together.
\end{abstract}

\section{The two sequences and the theorem}

\begin{definition}[Reversed binary word]
For $N\geq0$, let $\revbin(N)$ be the binary expansion of $N$ written from
least significant bit to most significant bit, represented as a word over
$0<1$.  Thus $\revbin(26)=01011$.
\end{definition}
\begin{leanbridge}
\begin{lstlisting}
abbrev Bit := Bool
def revBinary (n : Nat) : List Bit := n.bits
example : revBinary 26 =
  [false, true, false, true, true] := rfl
\end{lstlisting}
\end{leanbridge}

\begin{definition}[Lyndon word]
A nonempty word $w$ is Lyndon when it is lexicographically smaller than every
nontrivial cyclic rotation of itself.  Equivalently for the words used here,
$w$ is smaller than every proper nonempty suffix.
\end{definition}
\begin{leanbridge}
\begin{lstlisting}
def IsLyndon {α : Type*} [LinearOrder α] (w : List α) : Prop :=
  w ≠ [] ∧
    ∀ i, 0 < i → i < w.length → w < w.rotate i

/-- An auxiliary formulation that compares a word directly with each proper
nonempty suffix. For the binary lemmas below, this is often easier to use than
the rotation-based definition. -/
\end{lstlisting}
\end{leanbridge}

\begin{definition}[A328596]
Put
\[
  \SeqA=\{N\in\mathbb N: N>0\text{ and }\revbin(N)\text{ is Lyndon}\}.
\]
Write its increasing enumeration as $A(1)<A(2)<\cdots$.
\end{definition}
\begin{leanbridge}
\begin{lstlisting}
def InA328596 (n : Nat) : Prop :=
  0 < n /\ IsLyndon (revBinary n)
\end{lstlisting}
\end{leanbridge}

\begin{definition}[A348369]
For $n\geq1$, let $a(n)$ be the number of unordered pairs $\{k,m\}$ with
$0<k,m<n$ and
\[
                  A(n)=A(k)+A(m).
\]
This is OEIS A348369.  Its first term is $a(1)=0$.
\end{definition}
\begin{leanbridge}
The supplied Lean development formalizes the existence theorem
on the values $A(n)$ rather than the finite counting function $a(n)$.
The final corollary below explains why this is exactly what is needed for
positivity of the count.
\end{leanbridge}

\begin{theorem}[Main theorem]\label{thm:main}
For every $n>1$, one has $a(n)\geq1$.
\end{theorem}

\section{Binary structure}

\begin{lemma}[Powers of two]\label{lem:powers}
For every $k\geq0$,
\[
             \revbin(2^k)=0^k1,
\]
and $2^k\in\SeqA$.
\end{lemma}
\begin{proof}
Multiplication by $2$ prepends a zero to the least-significant-bit-first word.
Starting from $\revbin(1)=1$ gives the displayed form.  The word $0^k1$ is
smaller than each of its proper nonempty suffixes, hence is Lyndon.
\end{proof}
\begin{leanbridge}
\begin{lstlisting}
theorem revBinary_pow_two (k : Nat) :
    revBinary (2 ^ k) = List.replicate k false ++ [true] := by
  induction k with
  | zero =>
      simp [revBinary_one]
  | succ k ih =>
      have hk : 2 ^ k ≠ 0 := by
        have hk' : 0 < 2 ^ k := by
          positivity
        exact Nat.ne_of_gt hk'
      calc
        revBinary (2 ^ (k + 1)) = revBinary (2 * (2 ^ k)) := by
          rw [Nat.pow_succ, Nat.mul_comm]
        _ = false :: revBinary (2 ^ k) := by
          simp [revBinary_two_mul hk]
        _ = false :: (List.replicate k false ++ [true]) := by
          rw [ih]
        _ = List.replicate (k + 1) false ++ [true] := by
          simp [List.replicate]

/-- Multiplying a positive number by `2^k` prefixes `k` zero bits in the
reversed binary expansion. -/
\end{lstlisting}
\end{leanbridge}

\begin{lemma}[Odd-factor shape]\label{lem:shape}
If $N=2^k(2r+1)$, then
\[
 \revbin(N)=0^k1\,\revbin(r).
\]
If $r>0$ and the least-significant set bit $2^k$ is removed, then
\[
 \revbin(N-2^k)=0^{k+1}\revbin(r).
\]
\end{lemma}
\begin{proof}
The factor $2^k$ contributes $k$ initial zeros.  Since $2r+1$ is odd, its
first reversed bit is $1$ and the remaining bits are those of $r$.
Moreover $N-2^k=2^{k+1}r$, which gives the second identity.
\end{proof}
\begin{leanbridge}
\begin{lstlisting}
theorem revBinary_pow_two_mul_odd (k r : Nat) :
    revBinary (2 ^ k * (2 * r + 1)) = List.replicate k false ++ (true :: revBinary r) := by
  calc
    revBinary (2 ^ k * (2 * r + 1)) = List.replicate k false ++ revBinary (2 * r + 1) := by
      exact revBinary_mul_pow_two k (2 * r + 1) (by omega)
    _ = List.replicate k false ++ (true :: revBinary r) := by
      simp [revBinary_two_mul_add_one]

/-- Clearing the least-significant set bit of `2^k * (2r + 1)` for `r > 0`
adds one more leading `0` in reversed binary. -/
\end{lstlisting}
\end{leanbridge}

\section{The Lyndon transformation}

\begin{lemma}[Zero-block step]\label{lem:step}
Let $u$ be a binary word.  If $0^k1u1$ is Lyndon, then $0^{k+1}u1$ is
Lyndon.
\end{lemma}
\begin{proof}
Use the suffix characterization.  First delete the distinguished $1$ after
the initial zero block.  For a suffix beginning inside the zero block,
lexicographic comparison is determined after cancelling the common zeros;
the additional nonempty block of zeros is smaller than the terminal word
$u1$.  For a suffix beginning to the right of the distinguished $1$, the
required inequality follows transitively from the corresponding suffix
inequality for $0^k1u1$.  The resulting word remains suffix-Lyndon.  Prefixing
one additional $0$ preserves that property because every word considered
ends in $1$.  The equivalence of suffix and rotation formulations completes
the proof.
\end{proof}
\begin{leanbridge}
\begin{lstlisting}
theorem remove_first_true_preserves_isSuffixLyndon (k : Nat) (u : List Bool)
    (h : IsSuffixLyndon (List.replicate k false ++ (true :: (u ++ [true])))) :
    IsSuffixLyndon (List.replicate k false ++ (u ++ [true])) := by
  let t := u ++ [true]
  have hyx : List.replicate k false ++ t < List.replicate k false ++ (true :: t) := by
    simpa [t] using append_left_lt_of_lt (List.replicate k false) (append_true_lt_true_cons u)
  constructor
  · simp [t]
  · intro i hi0 hi
    by_cases hik : i ≤ k
    · have hdecomp : k = (k - i) + i := by omega
      have hy : List.replicate k false ++ t =
          List.replicate (k - i) false ++ (List.replicate i false ++ t) := by
        rw [hdecomp, List.replicate_add, List.append_assoc]
        simp
      have hdropy : (List.replicate k false ++ t).drop i = List.replicate (k - i) false ++ t := by
        rw [List.drop_append_of_le_length]
        · simp [hik, t]
        · simpa using hik
      calc
        List.replicate k false ++ t
            = List.replicate (k - i) false ++ (List.replicate i false ++ t) := hy
        _ < List.replicate (k - i) false ++ t :=
          append_left_lt_of_lt (List.replicate (k - i) false)
            (replicate_falses_lt_self_append_true_of_pos hi0 u)
        _ = (List.replicate k false ++ t).drop i := hdropy.symm
    · have hklt : k < i := lt_of_not_ge hik
      let j := i - k
      have hjpos : 0 < j := by
        dsimp [j]
        omega
      have hjlen : j < t.length := by
        dsimp [j] at *
        simp [t] at hi ⊢
        omega
      have hdropy : (List.replicate k false ++ t).drop i = t.drop j := by
        calc
          (List.replicate k false ++ t).drop i
              = ((List.replicate k false ++ t).drop k).drop j := by
                  rw [show i = k + j by
                    dsimp [j]
                    omega, List.drop_drop]
          _ = t.drop j := by
            rw [List.drop_append_of_le_length]
            · simp [t]
            · simp
      have hdropx : (List.replicate k false ++ (true :: t)).drop (i + 1) = t.drop j := by
        calc
          (List.replicate k false ++ (true :: t)).drop (i + 1)
              = ((List.replicate k false ++ (true :: t)).drop (k + 1)).drop j := by
                  rw [show i + 1 = (k + 1) + j by
                    dsimp [j]
                    omega, List.drop_drop]
          _ = t.drop j := by
            have hxk1 : (List.replicate k false ++ (true :: t)).drop (k + 1) = t := by
              calc
                (List.replicate k false ++ (true :: t)).drop (k + 1)
                    = ((List.replicate k false ++ (true :: t)).drop k).drop 1 := by
                        rw [show k + 1 = k + 1 by omega, List.drop_drop]
                _ = (true :: t).drop 1 := by
                  rw [List.drop_append_of_le_length]
                  · simp
                  · simp
                _ = t := by simp
            rw [hxk1]
      have hxlt : List.replicate k false ++ (true :: t) < t.drop j := by
        have hxlen : i + 1 < (List.replicate k false ++ (true :: t)).length := by
          simp [t] at hi ⊢
          omega
        have htmp := h.lt_drop (Nat.succ_pos _) hxlen
        rw [hdropx] at htmp
        exact htmp
      rw [hdropy]
      exact lt_trans hyx hxlt

/-- Suffix-form version of the main combinatorial transformation:
from `0^k 1 u 1`, delete the first `1` after the initial zero block, then
prepend one more `0`. -/
\end{lstlisting}
\end{leanbridge}

\begin{proposition}[Closure under clearing the least-significant set bit]
\label{prop:clear}
Let $r>0$.  If $N=2^k(2r+1)$ belongs to $\SeqA$, then
$N-2^k=2^{k+1}r$ also belongs to $\SeqA$.
\end{proposition}
\begin{proof}
Because $r>0$, its reversed binary word ends in $1$; write
$\revbin(r)=u1$.  By Lemma~\ref{lem:shape},
$\revbin(N)=0^k1u1$.  This word is Lyndon by $N\in\SeqA$.
Lemma~\ref{lem:step} shows that $0^{k+1}u1$ is Lyndon.  Lemma~\ref{lem:shape}
identifies it with $\revbin(N-2^k)$, and the latter number is positive.
\end{proof}
\begin{leanbridge}
\begin{lstlisting}
theorem clear_lsb_shape_mem_a328596 (k r : Nat) (hr : r ≠ 0)
    (hA : InA328596 (2 ^ k * (2 * r + 1))) :
    InA328596 (2 ^ (k + 1) * r) := by
  constructor
  · exact Nat.mul_pos (by positivity) (Nat.pos_of_ne_zero hr)
  · rcases revBinary_pos_ends_with_true (Nat.pos_of_ne_zero hr) with ⟨u, hu⟩
    have hword : IsSuffixLyndon (List.replicate k false ++ (true :: (u ++ [true]))) := by
      have hly : IsLyndon (List.replicate k false ++ (true :: (u ++ [true]))) := by
        simpa [revBinary_pow_two_mul_odd, hu] using hA.2
      exact hly.toIsSuffixLyndon
    have hsmall : IsLyndon (List.replicate (k + 1) false ++ (u ++ [true])) :=
      zero_block_step_isLyndon k u hword
    simpa [revBinary_mul_pow_two (k + 1) r hr, hu, List.append_assoc] using hsmall
\end{lstlisting}
\end{leanbridge}

\section{Additive decomposition and positivity}

\begin{theorem}[Additive decomposition]\label{thm:additive}
If $N>1$ and $N\in\SeqA$, then there exist $X,Y\in\SeqA$ such that
$N=X+Y$.
\end{theorem}
\begin{proof}
Write $N=2^kq$, where $q$ is odd.

If $q=1$, then $N=2^k$.  Since $N>1$, $k=j+1$ for some $j$, and
\[
                 N=2^j+2^j.
\]
Both summands belong to $\SeqA$ by Lemma~\ref{lem:powers}.

Now suppose $q>1$.  Write $q=2r+1$.  Then $r>0$.  The least-significant set
bit $X=2^k$ belongs to $\SeqA$ by Lemma~\ref{lem:powers}.  Proposition
\ref{prop:clear} shows that
$Y=2^{k+1}r=N-2^k$ also belongs to $\SeqA$.  Hence
\[
                 N=2^k+2^{k+1}r=X+Y,
\]
which proves the assertion.
\end{proof}
\begin{leanbridge}
\begin{lstlisting}
theorem a328596_additive {n : Nat} (hn : 1 < n) (hA : InA328596 n) :
    ∃ a b, InA328596 a ∧ InA328596 b ∧ n = a + b := by
  let k := padicValNat 2 n
  let q := n.divMaxPow 2
  have hfact : 2 ^ k * q = n := by
    simpa [k, q] using (pow_padicValNat_mul_divMaxPow 2 n)
  by_cases hq1 : q = 1
  · have hkshape : n = 2 ^ k := by
      simpa [hq1] using hfact.symm
    rcases Nat.eq_zero_or_pos k with hk0 | hkpos
    · have hnone : n = 1 := by simpa [hk0] using hkshape
      omega
    · obtain ⟨j, hj⟩ := Nat.exists_eq_succ_of_ne_zero (Nat.ne_of_gt hkpos)
      refine ⟨2 ^ j, 2 ^ j, two_pow_mem_a328596 j, two_pow_mem_a328596 j, ?_⟩
      calc
        n = 2 ^ (j + 1) := by simpa [hj] using hkshape
        _ = 2 ^ j + 2 ^ j := by
          rw [Nat.pow_succ]
          nth_rw 2 [show 2 = 1 + 1 by rfl]
          rw [Nat.mul_add]
          simp
  · have hqpos : 0 < q := by
      apply Nat.pos_of_ne_zero
      intro hq0
      rw [hq0, mul_zero] at hfact
      exact (Nat.ne_of_gt hA.1) hfact.symm
    have hqodd : ¬ 2 ∣ q := by
      simpa [q] using (Nat.not_dvd_divMaxPow (by decide : 1 < 2) (Nat.ne_of_gt hA.1))
    have hbodd : q.bodd = true := by
      cases hqbod : q.bodd <;> simp [Nat.mod_two_of_bodd, hqbod, Nat.dvd_iff_mod_eq_zero] at hqodd ⊢
    let r := q.div2
    have hqshape : q = 2 * r + 1 := by
      have hbit : Nat.bit true r = q := by
        simpa [r, hbodd] using (Nat.bit_bodd_div2 q)
      simpa [Nat.bit, r] using hbit.symm
    have hr : r ≠ 0 := by
      intro hr0
      have : q = 1 := by simpa [hqshape, r, hr0]
      exact hq1 this
    have hshape : n = 2 ^ k * (2 * r + 1) := by
      rw [← hfact, hqshape]
    have hsmall : InA328596 (2 ^ (k + 1) * r) := by
      have hAshape : InA328596 (2 ^ k * (2 * r + 1)) := by
        simpa [hshape] using hA
      exact clear_lsb_shape_mem_a328596 k r hr hAshape
    refine ⟨2 ^ k, 2 ^ (k + 1) * r, two_pow_mem_a328596 k, hsmall, ?_⟩
    calc
      n = 2 ^ k * (2 * r + 1) := hshape
      _ = 2 ^ k + 2 ^ (k + 1) * r := by
        rw [Nat.mul_add, Nat.mul_one, Nat.pow_succ]
        ac_rfl
\end{lstlisting}
\end{leanbridge}

\begin{proof}[Proof of Theorem~\ref{thm:main}]
Fix $n>1$ and put $N=A(n)$.  Since the increasing sequence $A$ begins with
$A(1)=1$, we have $N>1$.  Theorem~\ref{thm:additive} supplies
$X,Y\in\SeqA$ with $N=X+Y$.  Positivity gives $0<X,Y<N$.  Hence there are
indices $0<k,m<n$ with $X=A(k)$ and $Y=A(m)$.  Thus the unordered pair
$\{k,m\}$ is counted by $a(n)$, and therefore $a(n)\geq1$.
\end{proof}
\begin{leanbridge}
The machine-checked mathematical core is exactly
\texttt{a328596\_additive}.  A fully formal theorem mentioning
\texttt{a(n)} additionally requires Lean definitions of the increasing
enumeration of \texttt{InA328596} and of the finite quotient count of
unordered pairs.  Those bookkeeping definitions are not presented here.
\end{leanbridge}

\begin{remark}[Strength of the result]
The proof constructs a witness.  For non-powers of two it selects the
least-significant set bit $2^{v_2(N)}$ and the complementary summand.
Accordingly, it proves existence rather than merely deriving positivity from
computed terms.
\end{remark}
\begin{leanbridge}
The witness is exposed by the two branches of
\texttt{a328596\_additive}: either $(2^j,2^j)$ or
$(2^k,2^{k+1}r)$.
\end{leanbridge}

\section*{Reproducibility note}
The formal development consists of \texttt{BinaryWords.lean},
\texttt{Lyndon.lean}, and \texttt{Sequence.lean}, using Lean 4 and Mathlib.
Every nontrivial Lean declaration cited in the article is reproduced in full below.
For complete reproducibility, Appendix~\ref{app:lean-sources} also contains the
verbatim contents of all three Lean source files; no proof is abbreviated by an ellipsis.


\clearpage
\appendix
\section{Complete Lean 4 source files}\label{app:lean-sources}
This appendix reproduces the complete machine-checked development verbatim, including
imports, namespaces, definitions, auxiliary lemmas, sanity checks, and the final theorem.
The three modules are shown in dependency order.

\subsection{\texttt{BinaryWords.lean}}
\begin{lstlisting}
import Mathlib.Data.Nat.Bits
import Mathlib.Tactic

/-!
# Reversed Binary Words

This file defines the reversed binary expansion used in the OEIS A328596
formalization and proves a first structural lemma for powers of two.
-/

namespace A328596

/-- We represent binary words as lists of bits.
`false` means `0` and `true` means `1`. -/
abbrev Bit := Bool

/-- `revBinary n` is the ordinary binary expansion of `n`, written from
least significant bit to most significant bit.

For example, `26 = 11010₂`, so `revBinary 26 = [0, 1, 0, 1, 1]`. -/
def revBinary (n : Nat) : List Bit :=
  n.bits

@[simp] theorem revBinary_zero : revBinary 0 = [] := by
  rfl

@[simp] theorem revBinary_one : revBinary 1 = [true] := by
  simp [revBinary]

@[simp] theorem revBinary_two_mul {n : Nat} (hn : n ≠ 0) :
    revBinary (2 * n) = false :: revBinary n := by
  simp [revBinary, Nat.bit0_bits, hn]

@[simp] theorem revBinary_two_mul_add_one (n : Nat) :
    revBinary (2 * n + 1) = true :: revBinary n := by
  simp [revBinary]

/-- The reversed binary word of `2^k` is `0^k 1`. -/
theorem revBinary_pow_two (k : Nat) :
    revBinary (2 ^ k) = List.replicate k false ++ [true] := by
  induction k with
  | zero =>
      simp [revBinary_one]
  | succ k ih =>
      have hk : 2 ^ k ≠ 0 := by
        have hk' : 0 < 2 ^ k := by
          positivity
        exact Nat.ne_of_gt hk'
      calc
        revBinary (2 ^ (k + 1)) = revBinary (2 * (2 ^ k)) := by
          rw [Nat.pow_succ, Nat.mul_comm]
        _ = false :: revBinary (2 ^ k) := by
          simp [revBinary_two_mul hk]
        _ = false :: (List.replicate k false ++ [true]) := by
          rw [ih]
        _ = List.replicate (k + 1) false ++ [true] := by
          simp [List.replicate]

/-- Multiplying a positive number by `2^k` prefixes `k` zero bits in the
reversed binary expansion. -/
theorem revBinary_mul_pow_two (k n : Nat) (hn : n ≠ 0) :
    revBinary (2 ^ k * n) = List.replicate k false ++ revBinary n := by
  induction k with
  | zero =>
      simp
  | succ k ih =>
      have hk : 2 ^ k * n ≠ 0 := by
        exact Nat.mul_ne_zero (pow_ne_zero k (by decide)) hn
      have hmul : 2 ^ (k + 1) * n = 2 * (2 ^ k * n) := by
        rw [Nat.pow_succ]
        ac_rfl
      calc
        revBinary (2 ^ (k + 1) * n) = revBinary (2 * (2 ^ k * n)) := by rw [hmul]
        _ = false :: revBinary (2 ^ k * n) := by
          simp [revBinary_two_mul hk]
        _ = false :: (List.replicate k false ++ revBinary n) := by
          rw [ih]
        _ = List.replicate (k + 1) false ++ revBinary n := by
          simp [List.replicate]

/-- Splitting off the least-significant `1` from an odd factor exposes the
shape used in the additive decomposition argument. -/
theorem revBinary_pow_two_mul_odd (k r : Nat) :
    revBinary (2 ^ k * (2 * r + 1)) = List.replicate k false ++ (true :: revBinary r) := by
  calc
    revBinary (2 ^ k * (2 * r + 1)) = List.replicate k false ++ revBinary (2 * r + 1) := by
      exact revBinary_mul_pow_two k (2 * r + 1) (by omega)
    _ = List.replicate k false ++ (true :: revBinary r) := by
      simp [revBinary_two_mul_add_one]

/-- Clearing the least-significant set bit of `2^k * (2r + 1)` for `r > 0`
adds one more leading `0` in reversed binary. -/
theorem revBinary_clear_lsb (k r : Nat) (hr : r ≠ 0) :
    revBinary (2 ^ k * (2 * r + 1) - 2 ^ k) =
      List.replicate (k + 1) false ++ revBinary r := by
  have hsub : 2 ^ k * (2 * r + 1) - 2 ^ k = 2 ^ k * (2 * r) := by
    calc
      2 ^ k * (2 * r + 1) - 2 ^ k = 2 ^ k * (2 * r + 1) - 2 ^ k * 1 := by simp
      _ = 2 ^ k * ((2 * r + 1) - 1) := by rw [Nat.mul_sub_left_distrib]
      _ = 2 ^ k * (2 * r) := by simp
  have hmul : 2 ^ k * (2 * r) = 2 ^ (k + 1) * r := by
    rw [Nat.pow_succ]
    ac_rfl
  calc
    revBinary (2 ^ k * (2 * r + 1) - 2 ^ k)
        = revBinary (2 ^ k * (2 * r)) := by rw [hsub]
    _ = revBinary (2 ^ (k + 1) * r) := by rw [hmul]
    _ = List.replicate (k + 1) false ++ revBinary r := by
          exact revBinary_mul_pow_two (k + 1) r hr

/-- Every positive reversed binary expansion ends with a `1`. -/
theorem revBinary_pos_ends_with_true {n : Nat} (hn : 0 < n) :
    ∃ u, revBinary n = u ++ [true] := by
  induction n using Nat.strong_induction_on with
  | h n ih =>
      have hn0 : n ≠ 0 := Nat.ne_of_gt hn
      by_cases hodd : n.bodd = true
      · cases hdiv : n.div2 with
        | zero =>
            use []
            have hn1 : n = 1 := by
              have hbit : Nat.bit true 0 = n := by
                simpa [hodd, hdiv] using Nat.bit_bodd_div2 n
              simpa [Nat.bit] using hbit.symm
            simp [hn1, revBinary_one]
        | succ m =>
            have hsmall : Nat.succ m < n := by
              simpa [hdiv] using Nat.binaryRec_decreasing hn0
            rcases ih (Nat.succ m) hsmall (Nat.succ_pos _) with ⟨u, hu⟩
            use true :: u
            have hbit : Nat.bit true (Nat.succ m) = n := by
              simpa [hodd, hdiv] using Nat.bit_bodd_div2 n
            have hnexpr : n = 2 * Nat.succ m + 1 := by
              simpa [Nat.bit] using hbit.symm
            rw [hnexpr, revBinary_two_mul_add_one, hu]
            simp
      · have hodd0 : n.bodd = false := by
          cases hbod : n.bodd <;> simp_all
        cases hdiv : n.div2 with
        | zero =>
            have : False := by
              have hbit : Nat.bit false 0 = n := by
                simpa [hodd0, hdiv] using Nat.bit_bodd_div2 n
              have hnzero : n = 0 := by
                simpa [Nat.bit] using hbit.symm
              exact hn0 hnzero
            exact this.elim
        | succ m =>
            have hsmall : Nat.succ m < n := by
              simpa [hdiv] using Nat.binaryRec_decreasing hn0
            rcases ih (Nat.succ m) hsmall (Nat.succ_pos _) with ⟨u, hu⟩
            use false :: u
            have hbit : Nat.bit false (Nat.succ m) = n := by
              simpa [hodd0, hdiv] using Nat.bit_bodd_div2 n
            have hnexpr : n = 2 * Nat.succ m := by
              simpa [Nat.bit] using hbit.symm
            rw [hnexpr, revBinary_two_mul (by simp [hdiv]), hu]
            simp

example : revBinary 26 = [false, true, false, true, true] := rfl

end A328596
\end{lstlisting}

\subsection{\texttt{Lyndon.lean}}
\begin{lstlisting}
import Mathlib.Data.List.Lex
import Mathlib.Data.List.Infix
import Mathlib.Data.List.Rotate
import Mathlib.Tactic

/-!
# Lyndon Words

This file contains the transparent list-based definition of a Lyndon word used
throughout the project.
-/

namespace A328596

/-- A word is Lyndon if it is nonempty and is lexicographically strictly smaller
than every nontrivial cyclic rotation.

This is the direct list-based version of the mathematical definition. -/
def IsLyndon {α : Type*} [LinearOrder α] (w : List α) : Prop :=
  w ≠ [] ∧
    ∀ i, 0 < i → i < w.length → w < w.rotate i

/-- An auxiliary formulation that compares a word directly with each proper
nonempty suffix. For the binary lemmas below, this is often easier to use than
the rotation-based definition. -/
def IsSuffixLyndon {α : Type*} [LinearOrder α] (w : List α) : Prop :=
  w ≠ [] ∧
    ∀ i, 0 < i → i < w.length → w < w.drop i

/-- A Boolean checker for the Lyndon condition, used for small computational
sanity checks.

It checks the nonempty condition and then tests all rotations with cut positions
`i < w.length`, treating `i = 0` as the trivial rotation that we ignore. -/
def isLyndonBool {α : Type*} [LinearOrder α] (w : List α) : Bool :=
  decide (w ≠ []) &&
    (List.range w.length).all fun i => decide (i = 0 ∨ w < w.rotate i)

theorem IsLyndon.ne_nil {α : Type*} [LinearOrder α] {w : List α} (h : IsLyndon w) :
    w ≠ [] :=
  h.1

theorem IsLyndon.lt_rotate {α : Type*} [LinearOrder α] {w : List α} (h : IsLyndon w)
    {i : Nat} (hi0 : 0 < i) (hi : i < w.length) :
    w < w.rotate i :=
  h.2 i hi0 hi

theorem IsSuffixLyndon.ne_nil {α : Type*} [LinearOrder α] {w : List α}
    (h : IsSuffixLyndon w) : w ≠ [] :=
  h.1

theorem IsSuffixLyndon.lt_drop {α : Type*} [LinearOrder α] {w : List α}
    (h : IsSuffixLyndon w) {i : Nat} (hi0 : 0 < i) (hi : i < w.length) :
    w < w.drop i :=
  h.2 i hi0 hi

theorem cons_lt_cons {α : Type*} [LinearOrder α] {a : α} {u v : List α}
    (h : u < v) : a :: u < a :: v := by
  exact (List.lt_iff_lex_lt _ _).2 <| List.Lex.cons ((List.lt_iff_lex_lt _ _).1 h)

theorem append_right_lt_of_lt {α : Type*} [LinearOrder α] {u v : List α} (t : List α)
    (h : u < v) : u < v ++ t := by
  exact (List.lt_iff_lex_lt _ _).2 <|
    List.Lex.append_right (r := (· < ·)) t ((List.lt_iff_lex_lt _ _).1 h)

theorem append_left_lt_of_lt {α : Type*} [LinearOrder α] {u v : List α} (p : List α)
    (h : u < v) : p ++ u < p ++ v := by
  exact (List.lt_iff_lex_lt _ _).2 <|
    List.Lex.append_left (R := (· < ·)) ((List.lt_iff_lex_lt _ _).1 h) p

theorem append_left_cancel_lt {α : Type*} [LinearOrder α] {p u v : List α}
  (h : p ++ u < p ++ v) : u < v := by
  induction p with
  | nil => simpa using h
  | cons a p ih =>
    have h' : p ++ u < p ++ v := by
      exact (List.lex_cons_iff).1 h
    exact ih h'

theorem append_right_lt_of_lt_of_eq_length {α : Type*} [LinearOrder α]
  {u v : List α} (t : List α) (h : u < v) (hlen : u.length = v.length) :
  u ++ t < v ++ t := by
  induction u generalizing v with
  | nil =>
    cases v with
    | nil => cases lt_irrefl [] h
    | cons b v => simp at hlen
  | cons a u ih =>
    cases v with
    | nil => simp at hlen
    | cons b v =>
      rcases lt_or_eq_of_le (List.head_le_of_lt h) with hab | rfl
      · exact (List.lt_iff_lex_lt _ _).2 <| List.Lex.rel hab
      · have htail : u < v := by simpa [List.lex_cons_iff] using h
        have hlen' : u.length = v.length := by simpa using hlen
        exact cons_lt_cons (ih htail hlen')

theorem append_lt_of_lt_of_not_prefix {α : Type*} [LinearOrder α]
  {u v t : List α} (h : u < v) (hprefix : ¬ u <+: v) : u ++ t < v := by
  induction u generalizing v t with
  | nil =>
    simpa using hprefix
  | cons a u ih =>
    cases v with
    | nil => cases h
    | cons b v =>
      rcases lt_or_eq_of_le (List.head_le_of_lt h) with hab | rfl
      · exact (List.lt_iff_lex_lt _ _).2 <| List.Lex.rel hab
      · have htail : u < v := by simpa [List.lex_cons_iff] using h
        have hprefix' : ¬ u <+: v := by
          intro hp
          exact hprefix (List.cons_prefix_cons.mpr ⟨rfl, hp⟩)
        exact cons_lt_cons (ih htail hprefix')

theorem IsSuffixLyndon.toIsLyndon {α : Type*} [LinearOrder α] {w : List α}
    (h : IsSuffixLyndon w) : IsLyndon w := by
  constructor
  · exact h.1
  · intro i hi0 hi
    rw [List.rotate_eq_drop_append_take hi.le]
    exact append_right_lt_of_lt (List.take i w) (h.lt_drop hi0 hi)

theorem IsLyndon.drop_not_prefix {α : Type*} [LinearOrder α] {w : List α}
    (h : IsLyndon w) {i : Nat} (hi0 : 0 < i) (hi : i < w.length) :
    ¬ w.drop i <+: w := by
  intro hprefix
  let v := w.drop i
  let u := w.take i
  let z := w.drop v.length
  have hsplit : w = u ++ v := by
    simpa [u, v] using (List.take_append_drop i w).symm
  have hrot_i : w < v ++ u := by
    have hrot := h.lt_rotate hi0 hi
    rw [List.rotate_eq_drop_append_take hi.le] at hrot
    simpa [u, v] using hrot
  have hprefix_split : w = v ++ w.drop v.length := by
    exact List.prefix_append_drop hprefix
  have hzltu : z < u := by
    rw [hprefix_split] at hrot_i
    simpa [z] using append_left_cancel_lt hrot_i
  have hvpos : 0 < v.length := by
    simp [v]
    omega
  have hvlt : v.length < w.length := by
    simp [v]
    omega
  have hzlen : z.length = u.length := by
    simp [u, v, z]
    omega
  have hrot_v : w < z ++ v := by
    have hrot : v ++ z < z ++ v := by
      have hrot' := h.lt_rotate hvpos hvlt
      rw [hprefix_split, List.rotate_append_length_eq] at hrot'
      simpa [z] using hrot'
    exact hprefix_split.symm ▸ hrot
  have hback : z ++ v < w := by
    have htmp : z ++ v < u ++ v :=
      append_right_lt_of_lt_of_eq_length v hzltu hzlen
    exact hsplit.symm ▸ htmp
  exact (lt_asymm hrot_v hback).elim

theorem IsLyndon.toIsSuffixLyndon {α : Type*} [LinearOrder α] {w : List α}
    (h : IsLyndon w) : IsSuffixLyndon w := by
  constructor
  · exact h.1
  · intro i hi0 hi
    let v := w.drop i
    let u := w.take i
    have hrot : w < v ++ u := by
      have hrot' := h.lt_rotate hi0 hi
      rw [List.rotate_eq_drop_append_take hi.le] at hrot'
      simpa [u, v] using hrot'
    by_contra hnv
    have hvne : v ≠ w := by
      intro hEq
      have : v.length < w.length := by
        simp [v]
        omega
      simpa [hEq] using this
    have hvw : v < w := lt_of_le_of_ne (le_of_not_gt hnv) hvne
    by_cases hprefix : v <+: w
    · exact (h.drop_not_prefix hi0 hi) hprefix
    · have hcontra : v ++ u < w := append_lt_of_lt_of_not_prefix hvw hprefix
      exact (lt_asymm hrot hcontra).elim

theorem isLyndon_singleton {α : Type*} [LinearOrder α] (a : α) :
    IsLyndon [a] := by
  constructor
  · simp
  · intro i hi0 hi
    have hlt : i < 1 := by simpa using hi
    omega

theorem isSuffixLyndon_singleton {α : Type*} [LinearOrder α] (a : α) :
    IsSuffixLyndon [a] := by
  constructor
  · simp
  · intro i hi0 hi
    have hlt : i < 1 := by simpa using hi
    omega

/-- If a binary suffix-Lyndon word ends in `1`, then prefixing one more `0`
produces another suffix-Lyndon word.

This is one of the structural closure properties behind the least-significant-
bit decomposition. -/
theorem prepend_false_preserves_isSuffixLyndon (u : List Bool)
    (h : IsSuffixLyndon (u ++ [true])) :
    IsSuffixLyndon (false :: (u ++ [true])) := by
  constructor
  · simp
  · intro i hi0 hi
    cases i with
    | zero => cases Nat.lt_asymm hi0 hi0
    | succ j =>
        have hj : j < (u ++ [true]).length := by
          simpa using hi
        have hjle : j ≤ u.length := by
          simpa using hj
        have hdrop : (u ++ [true]).drop j = u.drop j ++ [true] := by
          simpa using List.drop_append_of_le_length (l₁ := u) (l₂ := [true]) hjle
        cases huj : u.drop j with
        | nil =>
            rw [show (false :: (u ++ [true])).drop (j + 1) = [true] by simp [hdrop, huj]]
            exact (List.lt_iff_lex_lt _ _).2 <| List.Lex.rel (show false < true by decide)
        | cons b t =>
            cases b with
            | false =>
                have hju : j < u.length := by
                  have hlen' : 0 < (u.drop j).length := by
                    simp [huj]
                  rw [List.length_drop] at hlen'
                  omega
                have htail : (u ++ [true]).drop (j + 1) = t ++ [true] := by
                  rw [← List.drop_drop, hdrop, huj]
                  simp
                have huw : u ++ [true] < t ++ [true] := by
                  have huw' : u ++ [true] < (u ++ [true]).drop (j + 1) :=
                    h.lt_drop (Nat.succ_pos _) (by simpa using Nat.succ_lt_succ hju)
                  rw [htail] at huw'
                  exact huw'
                rw [show (false :: (u ++ [true])).drop (j + 1) = false :: (t ++ [true]) by
                  simp [hdrop, huj]]
                exact cons_lt_cons huw
            | true =>
                rw [show (false :: (u ++ [true])).drop (j + 1) = true :: (t ++ [true]) by
                  simp [hdrop, huj]]
                exact (List.lt_iff_lex_lt _ _).2 <| List.Lex.rel (show false < true by decide)

theorem prepend_false_preserves_isLyndon (u : List Bool)
    (h : IsSuffixLyndon (u ++ [true])) :
    IsLyndon (false :: (u ++ [true])) :=
  (prepend_false_preserves_isSuffixLyndon u h).toIsLyndon

theorem append_true_lt_true_cons (u : List Bool) :
  u ++ [true] < true :: (u ++ [true]) := by
  induction u with
  | nil =>
    exact (List.lt_iff_lex_lt _ _).2 <| List.Lex.cons List.Lex.nil
  | cons b u ih =>
    cases b with
    | false =>
      exact (List.lt_iff_lex_lt _ _).2 <| List.Lex.rel (show false < true by decide)
    | true =>
      simpa using cons_lt_cons (a := true) ih

theorem prepend_false_lt_self_append_true (u : List Bool) :
  false :: (u ++ [true]) < u ++ [true] := by
  induction u with
  | nil =>
    exact (List.lt_iff_lex_lt _ _).2 <| List.Lex.rel (show false < true by decide)
  | cons b u ih =>
    cases b with
    | false =>
      simpa using cons_lt_cons (a := false) ih
    | true =>
      exact (List.lt_iff_lex_lt _ _).2 <| List.Lex.rel (show false < true by decide)

theorem replicate_falses_lt_self_append_true (k : Nat) (u : List Bool) :
  List.replicate (k + 1) false ++ (u ++ [true]) < u ++ [true] := by
  induction k with
  | zero =>
      simpa using prepend_false_lt_self_append_true u
  | succ k ih =>
      have h₁ : false :: (List.replicate (k + 1) false ++ (u ++ [true])) < false :: (u ++ [true]) :=
        cons_lt_cons (a := false) ih
      have h₂ : List.replicate (k + 2) false ++ (u ++ [true]) < false :: (u ++ [true]) := by
        simpa [List.replicate, List.append_assoc] using h₁
      exact lt_trans h₂ (prepend_false_lt_self_append_true u)

theorem replicate_falses_lt_self_append_true_of_pos {i : Nat} (hi : 0 < i) (u : List Bool) :
    List.replicate i false ++ (u ++ [true]) < u ++ [true] := by
  cases i with
  | zero => cases Nat.lt_asymm hi hi
  | succ j =>
      simpa using replicate_falses_lt_self_append_true j u

theorem remove_first_true_preserves_isSuffixLyndon (k : Nat) (u : List Bool)
    (h : IsSuffixLyndon (List.replicate k false ++ (true :: (u ++ [true])))) :
    IsSuffixLyndon (List.replicate k false ++ (u ++ [true])) := by
  let t := u ++ [true]
  have hyx : List.replicate k false ++ t < List.replicate k false ++ (true :: t) := by
    simpa [t] using append_left_lt_of_lt (List.replicate k false) (append_true_lt_true_cons u)
  constructor
  · simp [t]
  · intro i hi0 hi
    by_cases hik : i ≤ k
    · have hdecomp : k = (k - i) + i := by omega
      have hy : List.replicate k false ++ t =
          List.replicate (k - i) false ++ (List.replicate i false ++ t) := by
        rw [hdecomp, List.replicate_add, List.append_assoc]
        simp
      have hdropy : (List.replicate k false ++ t).drop i = List.replicate (k - i) false ++ t := by
        rw [List.drop_append_of_le_length]
        · simp [hik, t]
        · simpa using hik
      calc
        List.replicate k false ++ t
            = List.replicate (k - i) false ++ (List.replicate i false ++ t) := hy
        _ < List.replicate (k - i) false ++ t :=
          append_left_lt_of_lt (List.replicate (k - i) false)
            (replicate_falses_lt_self_append_true_of_pos hi0 u)
        _ = (List.replicate k false ++ t).drop i := hdropy.symm
    · have hklt : k < i := lt_of_not_ge hik
      let j := i - k
      have hjpos : 0 < j := by
        dsimp [j]
        omega
      have hjlen : j < t.length := by
        dsimp [j] at *
        simp [t] at hi ⊢
        omega
      have hdropy : (List.replicate k false ++ t).drop i = t.drop j := by
        calc
          (List.replicate k false ++ t).drop i
              = ((List.replicate k false ++ t).drop k).drop j := by
                  rw [show i = k + j by
                    dsimp [j]
                    omega, List.drop_drop]
          _ = t.drop j := by
            rw [List.drop_append_of_le_length]
            · simp [t]
            · simp
      have hdropx : (List.replicate k false ++ (true :: t)).drop (i + 1) = t.drop j := by
        calc
          (List.replicate k false ++ (true :: t)).drop (i + 1)
              = ((List.replicate k false ++ (true :: t)).drop (k + 1)).drop j := by
                  rw [show i + 1 = (k + 1) + j by
                    dsimp [j]
                    omega, List.drop_drop]
          _ = t.drop j := by
            have hxk1 : (List.replicate k false ++ (true :: t)).drop (k + 1) = t := by
              calc
                (List.replicate k false ++ (true :: t)).drop (k + 1)
                    = ((List.replicate k false ++ (true :: t)).drop k).drop 1 := by
                        rw [show k + 1 = k + 1 by omega, List.drop_drop]
                _ = (true :: t).drop 1 := by
                  rw [List.drop_append_of_le_length]
                  · simp
                  · simp
                _ = t := by simp
            rw [hxk1]
      have hxlt : List.replicate k false ++ (true :: t) < t.drop j := by
        have hxlen : i + 1 < (List.replicate k false ++ (true :: t)).length := by
          simp [t] at hi ⊢
          omega
        have htmp := h.lt_drop (Nat.succ_pos _) hxlen
        rw [hdropx] at htmp
        exact htmp
      rw [hdropy]
      exact lt_trans hyx hxlt

/-- Suffix-form version of the main combinatorial transformation:
from `0^k 1 u 1`, delete the first `1` after the initial zero block, then
prepend one more `0`. -/
theorem zero_block_step_isSuffixLyndon (k : Nat) (u : List Bool)
    (h : IsSuffixLyndon (List.replicate k false ++ (true :: (u ++ [true])))) :
    IsSuffixLyndon (List.replicate (k + 1) false ++ (u ++ [true])) := by
  have h' : IsSuffixLyndon ((List.replicate k false ++ u) ++ [true]) := by
    simpa [List.append_assoc] using remove_first_true_preserves_isSuffixLyndon k u h
  simpa [List.replicate, List.replicate_add, List.append_assoc] using
    prepend_false_preserves_isSuffixLyndon (List.replicate k false ++ u) h'

/-- Rotation-form corollary of `zero_block_step_isSuffixLyndon`. -/
theorem zero_block_step_isLyndon (k : Nat) (u : List Bool)
    (h : IsSuffixLyndon (List.replicate k false ++ (true :: (u ++ [true])))) :
    IsLyndon (List.replicate (k + 1) false ++ (u ++ [true])) :=
  (zero_block_step_isSuffixLyndon k u h).toIsLyndon

theorem zeroes_then_one_isSuffixLyndon (k : Nat) :
    IsSuffixLyndon (List.replicate k false ++ [true]) := by
  induction k with
  | zero =>
      simpa using isSuffixLyndon_singleton true
  | succ k ih =>
      simpa [List.replicate, List.replicate_add] using
        prepend_false_preserves_isSuffixLyndon (List.replicate k false) ih

end A328596
\end{lstlisting}

\subsection{\texttt{Sequence.lean}}
\begin{lstlisting}
import Mathlib.Data.Nat.PadicValNat
import Mathlib.Tactic
import A328596.BinaryWords
import A328596.Lyndon

/-!
# The Sequence A328596

This file packages the reversed-binary Lyndon condition into a predicate on
natural numbers and records a first OEIS sanity check.
-/

namespace A328596

/-- Membership in OEIS A328596:
`n` is positive and its reversed binary expansion is a Lyndon word. -/
def InA328596 (n : Nat) : Prop :=
  0 < n ∧ IsLyndon (revBinary n)

/-- A Boolean version of `InA328596`, used for computed examples. -/
def inA328596Bool (n : Nat) : Bool :=
  decide (0 < n) && isLyndonBool (revBinary n)

/-- A small computed prefix, used only as a sanity check against OEIS. -/
def firstTermsUpTo (n : Nat) : List Nat :=
  (List.range (n + 1)).filter inA328596Bool

example : firstTermsUpTo 32 = [1, 2, 4, 6, 8, 12, 14, 16, 20, 24, 26, 28, 30, 32] := by
  decide

theorem zeroes_then_one_isLyndon (k : Nat) :
    IsLyndon (List.replicate k false ++ [true]) := by
  exact (zeroes_then_one_isSuffixLyndon k).toIsLyndon

theorem two_pow_mem_a328596 (k : Nat) : InA328596 (2 ^ k) := by
  constructor
  · have hk : 0 < 2 ^ k := by
      positivity
    exact hk
  · rw [revBinary_pow_two]
    exact zeroes_then_one_isLyndon k

theorem clear_lsb_shape_mem_a328596 (k r : Nat) (hr : r ≠ 0)
    (hA : InA328596 (2 ^ k * (2 * r + 1))) :
    InA328596 (2 ^ (k + 1) * r) := by
  constructor
  · exact Nat.mul_pos (by positivity) (Nat.pos_of_ne_zero hr)
  · rcases revBinary_pos_ends_with_true (Nat.pos_of_ne_zero hr) with ⟨u, hu⟩
    have hword : IsSuffixLyndon (List.replicate k false ++ (true :: (u ++ [true]))) := by
      have hly : IsLyndon (List.replicate k false ++ (true :: (u ++ [true]))) := by
        simpa [revBinary_pow_two_mul_odd, hu] using hA.2
      exact hly.toIsSuffixLyndon
    have hsmall : IsLyndon (List.replicate (k + 1) false ++ (u ++ [true])) :=
      zero_block_step_isLyndon k u hword
    simpa [revBinary_mul_pow_two (k + 1) r hr, hu, List.append_assoc] using hsmall

theorem a328596_additive {n : Nat} (hn : 1 < n) (hA : InA328596 n) :
    ∃ a b, InA328596 a ∧ InA328596 b ∧ n = a + b := by
  let k := padicValNat 2 n
  let q := n.divMaxPow 2
  have hfact : 2 ^ k * q = n := by
    simpa [k, q] using (pow_padicValNat_mul_divMaxPow 2 n)
  by_cases hq1 : q = 1
  · have hkshape : n = 2 ^ k := by
      simpa [hq1] using hfact.symm
    rcases Nat.eq_zero_or_pos k with hk0 | hkpos
    · have hnone : n = 1 := by simpa [hk0] using hkshape
      omega
    · obtain ⟨j, hj⟩ := Nat.exists_eq_succ_of_ne_zero (Nat.ne_of_gt hkpos)
      refine ⟨2 ^ j, 2 ^ j, two_pow_mem_a328596 j, two_pow_mem_a328596 j, ?_⟩
      calc
        n = 2 ^ (j + 1) := by simpa [hj] using hkshape
        _ = 2 ^ j + 2 ^ j := by
          rw [Nat.pow_succ]
          nth_rw 2 [show 2 = 1 + 1 by rfl]
          rw [Nat.mul_add]
          simp
  · have hqpos : 0 < q := by
      apply Nat.pos_of_ne_zero
      intro hq0
      rw [hq0, mul_zero] at hfact
      exact (Nat.ne_of_gt hA.1) hfact.symm
    have hqodd : ¬ 2 ∣ q := by
      simpa [q] using (Nat.not_dvd_divMaxPow (by decide : 1 < 2) (Nat.ne_of_gt hA.1))
    have hbodd : q.bodd = true := by
      cases hqbod : q.bodd <;> simp [Nat.mod_two_of_bodd, hqbod, Nat.dvd_iff_mod_eq_zero] at hqodd ⊢
    let r := q.div2
    have hqshape : q = 2 * r + 1 := by
      have hbit : Nat.bit true r = q := by
        simpa [r, hbodd] using (Nat.bit_bodd_div2 q)
      simpa [Nat.bit, r] using hbit.symm
    have hr : r ≠ 0 := by
      intro hr0
      have : q = 1 := by simpa [hqshape, r, hr0]
      exact hq1 this
    have hshape : n = 2 ^ k * (2 * r + 1) := by
      rw [← hfact, hqshape]
    have hsmall : InA328596 (2 ^ (k + 1) * r) := by
      have hAshape : InA328596 (2 ^ k * (2 * r + 1)) := by
        simpa [hshape] using hA
      exact clear_lsb_shape_mem_a328596 k r hr hAshape
    refine ⟨2 ^ k, 2 ^ (k + 1) * r, two_pow_mem_a328596 k, hsmall, ?_⟩
    calc
      n = 2 ^ k * (2 * r + 1) := hshape
      _ = 2 ^ k + 2 ^ (k + 1) * r := by
        rw [Nat.mul_add, Nat.mul_one, Nat.pow_succ]
        ac_rfl

end A328596
\end{lstlisting}

\begin{thebibliography}{9}
\bibitem{A348369}
OEIS Foundation Inc., \emph{The On-Line Encyclopedia of Integer Sequences},
Sequence A348369, \url{https://oeis.org/A348369}.
\bibitem{A328596}
OEIS Foundation Inc., \emph{The On-Line Encyclopedia of Integer Sequences},
Sequence A328596, \url{https://oeis.org/A328596}.
\bibitem{lean}
The Lean Community, \emph{Lean 4 and Mathlib},
\url{https://leanprover-community.github.io/}.
\end{thebibliography}

\end{document}
